Term papers and theses are prepared in DIN~A4 format (\qtyproduct{210 x
297}{\milli\metre}). The paper format is already predefined by the class option
\texttt{a4paper} in the template (see Listing~\ref{lst:tpl:thesis}).

Exceptions apply to dissertations and post-doctoral theses that are
(have to) be published by a book publisher at the same time. So if you are not
writing a dissertation or habilitation, skip the rest of this section.

Book publications appear in a smaller format. The exact dimensions vary
depending on the publisher. The university publisher uses for series of
publications usually use the DIN~A5 format (\qtyproduct{210 x 297}{\milli\metre}).
Authors of such works now face a problem. Reviewers usually require the work in
DIN~A4 format, but the publication takes place in a different format. There are
two possible solutions, both of which have their advantages and disadvantages.

\begin{enumerate}
\item You create the paper for the reviewers in DIN~A4 format, and the
      the manuscript is converted to the book format. This means additional
      work, but avoids certain problems that may arise during printing.

\item You submit the work in DIN~A4 format to the printer, who scales it down
      to the correct format for you.
      You then do not have to revise your work again. However, it is of great
      importance to take this route into account when creating the A4 template
      (in other words from the beginning).
\end{enumerate}

If you decide to go the first way, we recommend that you create the leader for
for the A5 version as shown in Listing~\ref{lst:tpl:papersize_a5}. Before
printing, please agree with your publisher on the value of the binding
correction (see section~\ref{sec:tpl:bcor}).

\lstinputlisting[language={[LaTeX]{TeX}},style=block,float,caption={conversion to DIN~A5},label={lst:tpl:papersize_a5}]{code/tpl_papersize_a5.tex}

If you want to take the less time-consuming second route, select at least the
basic size for the A4~template at least the basic font size~\qty{12}{\pt}. When
scaling from A4 to A5, the document will be reduced to approx.~\qty{71}{\percent}
of its native size. A character of size~\qty{12}{\pt} is then only barely
\qty{9}{\pt} in size. However, \LaTeX\ uses a separate font for each font size.
A character reduced from \qty{12}{\pt} to \qty{9}{\pt} has a different
appearance than a character that is used for \qty{9}{\pt}, designed for
\qty{9}{\pt}. When scaling down, the lines are lines are often thinner than in
the respective design size. The font "`Latin Modern"' is particularly
susceptible to this (see section~\ref{sec:tpl:font}).

If the basic size in the DIN~A4 template is too small, the characters in the A5
version will be too thin, and thinner lines may even disappear altogether.
The publisher will certainly point out to you during the test print that you
are using fonts outside their design size and that they cannot guarantee
optimal and that they cannot guarantee optimal printing as a result. You must
accept this warning.
